\begin{afterword}
\summaryhead

The post argues that calling quantum physics ``weird'' is a mistake about where the fault lies. Reality was running long before humans evolved, and the author describes it as a cloud of complex amplitude in configuration space. What is odd is the human expectation that the world is made of little billiard balls, an intuition that evolution built for chasing prey among solid objects.

Since surprise is the mark of a poor model, and intuition is a model, the goal is to reshape one's intuitions until reality looks normal. This cannot be forced, and pretending that quantum physics feels natural only denies the confusion; but dwelling on how bizarre it is wastes effort and keeps one in the old frame. The same holds outside physics: if you ``just don't understand'' how a physicist can believe in astrology, your model of people is at fault. Surprise is in the map, not the territory. Once a fact has been checked, the problem is the model.

\responsehead

The post's core advice is sound, and it has a distinguished precedent. When a checked fact surprises you, the fault is in your model, and the task is to change the model rather than resent the fact. Feynman gave the same advice in \textsc{QED}: if you do not like how nature is, ``that's going to get in the way of your understanding it.'' The post's closing caveat, to check first whether the alleged fact is a fact, is the right one.

The opening verdict goes further than the advice. Whoever calls quantum physics weird, the author thinks, ``will never understand physics no matter how many books they read.'' Feynman, in the same passage, called nature ``absurd from the point of view of common sense.'' The word did not stop him from understanding the theory. What matters on his account is whether one accepts the theory, not whether one finds it strange. The post's own fifth paragraph agrees: it is pointless to pretend that quantum physics feels natural ``when in fact it feels strange,'' and pretending is ``denying your confusion''.

The post also states as fact a contested view of what quantum mechanics describes. Reality, it says, ``in fact'' is ``a perfectly normal cloud of complex amplitude in configuration space'', and at the fundamental level there are no persistent objects moving in three dimensions. Both statements depend on the interpretation. Physicists and philosophers have reached no consensus on it; on readings in which the quantum state describes our knowledge, it is not an element of reality, and in pilot-wave (Bohmian) theory particles have positions in ordinary space. The book argues for its own reading only later, in the sequence on quantum physics. Here the reading is presented as the thing a reader must stop calling weird.

The claims about the mind are asserted. That thinking ``How bizarre!'' wastes time, throws you back into the old frame and feeds indignation is plausible, but the post gives no evidence for it.

In short: sound advice to change your model when a checked fact surprises you, together with a verdict on those who call quantum physics weird that Feynman's case contradicts, and with one contested interpretation of quantum mechanics stated as fact.
\end{afterword}
